\bookchapter{Capítulo 30: Espelhamento de tela, apresentações e conectividade veicular}

\chapterintro{A tela do celular tem no máximo 7 polegadas, mas sua TV tem 55 polegadas e seu projetor tem 100 polegadas. “Jogar” a imagem do celular para uma tela grande proporciona uma experiência completamente diferente. Este capítulo explica com clareza três cenários frequentes — espelhamento de tela, apresentações e conectividade veicular — para estender a capacidade do seu celular a telas maiores.}

\booksection{Os três campos do espelhamento sem fio — não os confunda}

O espelhamento sem fio parece ser sempre “do celular para a TV”, mas há três tecnologias subjacentes diferentes, que não são interoperáveis:

\begin{center}
\begin{tabularx}{\linewidth}{|X|X|X|}
\hline
\textbf{Tecnologia} & \textbf{Quem usa} & \textbf{Características} \tabularnewline
\hline
AirPlay & iPhone/iPad/Mac & Ecossistema Apple, a experiência mais fluida. Requer Apple TV ou TV compatível com AirPlay \tabularnewline
Miracast & Celulares Android & Conexão direta via Wi-Fi Direct, não requer roteador. O "espelhamento sem fio" padrão do Android usa este \tabularnewline
DLNA & Todas as plataformas & Apenas transmite arquivos de mídia (vídeo/música/fotos), não espelha a tela. Ideal para enviar vídeos para a TV \tabularnewline
\hline
\end{tabularx}
\end{center}

\booksubsection{Métodos de operação por plataforma:}

\begin{itemize}
  \item \textbf{iPhone}: Central de Controle > Espelhar Tela > Selecione sua TV/Apple TV. Se a TV não for compatível com AirPlay, compre um Apple TV ou um dongle compatível com AirPlay.
  \item \textbf{Android}: Deslize para baixo a barra de notificações > Encontre o botão "Espelhamento sem fio"/"Smart View"/"Multi-tela" > Selecione a TV. É necessário que o celular e a TV estejam na mesma rede Wi-Fi.
  \item \textbf{Solução universal}: Se a TV for muito antiga e não suportar nenhum espelhamento sem fio, compre um Chromecast ou Xiaomi Box/Huawei Box e conecte na porta HDMI da TV para receber o espelhamento.
\end{itemize}

\booksection{Espelhamento com fio — mais estável, latência zero}

O espelhamento sem fio é conveniente, mas tem latência, o que prejudica a experiência em jogos e apresentações. O espelhamento com fio resolve esse problema:

\begin{itemize}
  \item \textbf{iPhone}: É necessário um cabo adaptador Lightning/USB-C para HDMI (cerca de 100-200 RMB). Conectado, a tela do celular é transmitida diretamente para a porta HDMI da TV/projetor.
  \item \textbf{Android}: Alguns celulares suportam conexão direta USB-C para HDMI (requer suporte a DP Alt Mode, a maioria dos flagships suporta). Basta comprar um cabo USB-C para HDMI.
  \item \textbf{Cenário de sala de reunião}: Se você frequentemente precisa fazer apresentações com o celular, é recomendável ter sempre um cabo adaptador HDMI na bolsa. As TVs/projetores das salas de reunião quase sempre têm porta HDMI — é só conectar e usar, muito mais confiável do que se estressar com espelhamento sem fio.
\end{itemize}

\booksection{Transforme o celular em computador — Modo Desktop}

Alguns celulares Android suportam o “Modo Desktop” — ao conectar a um monitor, o celular exibe uma interface semelhante à de um computador desktop, que pode ser operada com mouse e teclado:

\begin{itemize}
  \item \textbf{Samsung DeX}: Recurso exclusivo dos flagships Samsung. Ao conectar a um monitor, surge uma interface desktop com suporte a múltiplas janelas — a experiência mais próxima de um computador real. Com teclado e mouse Bluetooth, é possível viajar a trabalho sem levar notebook.
  \item \textbf{Modo Desktop Huawei}: As séries Huawei Mate/Pura também têm função semelhante ao conectar a um monitor. Suporta múltiplas janelas, gerenciamento de arquivos e apps de escritório.
  \item \textbf{Motorola Ready For}: Os flagships Moto também suportam, com funcionalidade semelhante ao DeX.
  \item \textbf{Cenários aplicáveis}: Trabalho leve (responder e-mails, editar documentos), apresentação de PowerPoint, ver vídeos em tela grande. Não é adequado para operações pesadas de Excel ou design profissional.
\end{itemize}

\booksection{Conectividade veicular — a maneira correta de usar o celular ao dirigir}

Operar o celular enquanto dirige é ilegal, mas navegação, música e chamadas são necessidades reais. Os sistemas de conectividade veicular permitem usar as funções essenciais do celular na tela do carro, mantendo o celular guardado, sem tocá-lo:

\begin{itemize}
  \item \textbf{Apple CarPlay}: A escolha principal para usuários de iPhone. Conecta-se via USB ou sem fio ao sistema do carro, e uma interface simplificada do iPhone aparece na tela do veículo — navegação (Amap/Baidu Maps), música, telefone, mensagens, podcasts. O controle por voz Siri permite operação totalmente mãos-livres.
  \item \textbf{Android Auto}: Solução correspondente para usuários Android. Funcionalidade semelhante ao CarPlay, compatível com Google Maps, Spotify, etc. Na China, algumas funções não estão disponíveis devido às restrições dos serviços Google.
  \item \textbf{Huawei HiCar}: Solução própria da Huawei, com alta taxa de compatibilidade em modelos de veículos chineses. Além de navegação e música, oferece recursos exclusivos como videochamadas Huawei e integração com câmeras de estacionamento.
  \item \textbf{Baidu CarLife}: Solução multiplataforma, compatível com iOS e Android. Tem menos funções, mas ampla compatibilidade — muitos modelos de carros mais antigos só suportam CarLife.
\end{itemize}

\begin{tcbwarn}{Alerta de segurança}
Independentemente da solução de conectividade veicular usada, sempre priorize o controle por voz ao dirigir. Defina o destino da navegação com antecedência, escolha a playlist de música antes de sair e não opere a tela do veículo enquanto estiver em movimento.

\end{tcbwarn}

\booksection{Espelhamento em sala de reunião na prática — sem mais pânico}

Reuniões são o cenário mais frequente de espelhamento e também o mais propenso a problemas. Aqui estão algumas dicas práticas:

\begin{itemize}
  \item \textbf{Chegue cedo e teste}: Nunca espere a reunião começar para conectar o espelhamento. Chegue 10 minutos antes e teste o método de conexão. A probabilidade de falha no espelhamento Wi-Fi é muito maior do que você imagina.
  \item \textbf{Tenha um Plano B com fio}: Leve um cabo adaptador HDMI na bolsa. Quando o espelhamento sem fio não funcionar, o cabo com fio será sua tábua de salvação.
  \item \textbf{Desative a prévia de notificações}: Durante o espelhamento, suas mensagens do WeChat aparecerão na tela grande da sala de reunião. Antes de espelhar, desative a "Prévia de conteúdo nas notificações da tela de bloqueio" nas configurações, ou ative o "Modo Apresentação" (alguns celulares suportam ocultar notificações automaticamente durante o espelhamento).
  \item \textbf{Use "apresentação de documento" em vez de "operar o celular"}: Se você vai mostrar um PowerPoint, salve o arquivo PPT diretamente no celular e abra em tela cheia com o WPS, em vez de espelhar toda a área de trabalho do celular. Isso é mais profissional e mais fluido.
\end{itemize}

\begin{tcbpractice}{Mãos à obra: Faça um espelhamento do celular para a TV}
\begin{enumerate}
  \item Certifique-se de que o celular e a TV estão na mesma rede Wi-Fi.
  \item Abra uma foto ou um vídeo no celular.
  \item Tente o espelhamento sem fio (iPhone use Espelhar Tela, Android use o botão de espelhamento sem fio).
  \item Se falhar, verifique se a TV é compatível com o protocolo correspondente ou tente usar DLNA (toque no botão "Transmitir" no app de fotos/vídeos).
  \item Após o espelhamento bem-sucedido, experimente fazer outras operações no celular e veja se a tela grande sincroniza em tempo real.
\end{enumerate}

\end{tcbpractice}

\begin{tcbsummary}{Resumo do capítulo}
\begin{itemize}
  \item Três campos do espelhamento sem fio: AirPlay (Apple), Miracast (Android), DLNA (transmissão de mídia multiplataforma)
  \item Espelhamento com fio é mais estável e sem latência — cabo HDMI é essencial para salas de reunião
  \item Samsung DeX/Modo Desktop Huawei transformam o celular em um computador de trabalho leve
  \item Conectividade veicular: quatro soluções principais — CarPlay/Android Auto/HiCar/CarLife
  \item Regras de ouro para espelhamento em reunião: testar antes, ter Plano B, desativar prévia de notificações
\end{itemize}

\end{tcbsummary}
